% ================================
% \section*{\centering \bf Part II: Spatial-aware Representation}
% \addcontentsline{toc}{section}{Part II: Spatial-aware Representation}
% ================================





\begin{figure}
    \centering
    % \includegraphics[width=\linewidth, height=12cm]{example-image-golden}
    % \includegraphics[width=0.99\linewidth]{figures/part2.png}
    \includegraphics[width=\linewidth]{figures/part2_opener_v2.pdf}
    
    \caption{Part II in one view. (1) Encode first: a scene must be encoded before it can be perceived, reasoned about, or acted on. (2) One modality at a time: the same object as pixels, depth, points, meshes, voxels, and Gaussians, and then as language and sound, from explicit geometry to implicit structure (Chapter~\ref{chapter:single_modality}). (3) Into a shared space: the separate encodings are aligned in one space without losing the geometry that makes them spatially meaningful (Chapter~\ref{chapter:cross_modality}).}
    \label{fig:p2_main}
\end{figure}


Everything a spatially intelligent system does, perceiving a scene, reasoning over it, generating or acting within it, passes through some representation of space.
That representation is upstream of the capability: what a system can perceive is bounded by what its encoding can express, and what it can reason about is bounded by what that encoding makes accessible.
Before the later parts turn to perception, reasoning, and generation, this part asks the prior question of how spatial structure is captured at all, how the geometry, topology, and semantics of an environment come to be held in a form a machine can compute with.

The question is approached in two movements.
The first takes each modality on its own terms, asking how images, video, 3D scenes, language, and sound each encode spatial structure and what each renders explicit or leaves implicit.
The second asks how these separate encodings are brought into a shared space, so that signals as different as a pixel grid and a sentence can be aligned without losing the geometry that makes them spatially meaningful.
The order is deliberate, since a modality must be understood by itself before its representation can be fused with others, and it is in that fusion that the partial, sensor-bound view offered by any one signal begins to reinforce the rest.
Together they form the representational substrate on which the remainder of the handbook is built.

% Building upon this, spatial representation
% learning focuses on constructing effective and generalizable spatial abstractions that capture geometric,
% topological, and semantic properties of the environment.

\newpage


\chapter{Single-Modal Spatial Representation}
\label{chapter:single_modality}


% \lettrine[lines=3]{\initfamily\textcolor{darkgreen}{S}}{ingle-modal} 
\FramedInitial{S}{ingle-modal} spatial representations ask how much of a scene's spatial structure can be recovered from one modality alone, without recourse to cross-modal correspondence.
The question matters because every multimodal system is ultimately built from unimodal encoders, and the spatial content a fused representation can carry is bounded by what each modality already captures on its own.

What separates the modalities is not whether they carry spatial information but how overtly they carry it.
Vision presents geometry almost directly, arranging appearance and shape over a spatial grid or the volume a scene occupies.
Language and sound hold the same structure only obliquely, the one as symbolic relations among named entities, the other as the physical trace that propagation leaves on a signal.
Others still are tied to a particular sensor or task, trading generality for fidelity to the data they were built around.
This chapter proceeds along that gradient, from the most explicit geometric encodings of space toward the most implicit, examining each modality on its own terms before the question of how they combine is taken up.





\section{2D}
\label{sec:rep_2d}

\FramedInitial{T}{he} 2D representations are the most direct and information-dense entry point into spatial intelligence: images and videos provide a regular, discretely sampled grid over the visual field, on top of which a family of structured maps can be defined to expose geometric and semantic properties of the underlying scene.
We organize this section around seven such representations (Figure~\ref{fig:p2_s4_1_2d_representations}): the raw \emph{image/video tensor} itself, and six derived maps: \emph{segmentation}, \emph{depth}, \emph{surface normal}, \emph{edge}, \emph{optical flow}, and \emph{2D keypoint} representations.
All seven share the image grid as their domain, which makes them straightforward to fuse with one another and to consume with convolutional or transformer backbones; what differs is what each representation \emph{exposes}: appearance, semantics, geometry, boundaries, motion, or sparse structure.



\begin{figure}[!t]
    \centering
    \includegraphics[width=0.90\linewidth]{figures/part2/2D_representations.pdf}
    \caption{The seven 2D representations of this section, all defined over the image grid $\Omega \subset \mathbb{R}^2$: image/video, segmentation, depth, surface normals, edges, optical flow, and keypoints.}
    \label{fig:p2_s4_1_2d_representations}
\end{figure}

\subsection{Image and Video Representations} 
\label{subsec:p-2-sec-4-2D-image-and-video-representations}

An image can be modeled as a discrete sampling of a continuous visual signal over a 2D spatial domain. Formally, a digital image is represented as a function
\begin{equation}
I: \Omega \rightarrow \mathbb{R}^C, \quad \Omega \subset \mathbb{R}^2,
\end{equation}
where $\Omega$ denotes the image grid, $C$ is the number of channels (\eg, $C=1$ for grayscale, $C=3$ for RGB), and each pixel value encodes intensity or color information at a discrete spatial location \cite{gonzalez2011,szeliski2022}. 
In practice, images are stored as finite tensors $I \in \mathbb{R}^{H \times W \times C}$, where $H$ and $W$ denote image height and width.


A video extends this representation by introducing a temporal dimension.
Given a sequence of frames sampled over time, a video can be written as
\begin{equation}
V: \Omega \times \mathcal{T} \rightarrow \mathbb{R}^C,
\end{equation}
or equivalently as a 4D tensor $V \in \mathbb{R}^{T \times H \times W \times C}$, where $T$ denotes the number of frames \cite{szeliski2022}. 
This formulation preserves spatial structure within each frame while enabling the modeling of temporal dynamics across frames.
From a computational perspective, images and videos are treated as structured arrays of numerical values that can be efficiently processed using linear algebra operations and convolutional kernels. 
This regular grid structure enables weight sharing, local receptive fields, and efficient parallel computation, forming the basis of CNNs and their spatiotemporal extensions~\cite{lecun1998gradient,karpathy2014large}. 
In modern deep learning systems, image and video tensors may undergo normalization, resizing, patch embedding, or tokenization before being fed into neural
architectures, including CNNs and vision transformers~\cite{vit, liu2022convnet}.
For video inputs, spatiotemporal extensions either inflate 2D kernels into 3D convolutions~\cite{tran2015learning, carreira2017quo} or extend self-attention across space and time, as in the family of video transformers~\cite{arnab2021vivit, liu2022videoswin}.






\subsection{Segmentation Maps}
\label{subsec:p-2-sec-4-2D-segmentation-maps}
A segmentation map is a semantic representation that assigns a label or a mask to each pixel, thereby partitioning visual data into meaningful regions such as objects, object parts, or background categories~\cite{szeliski2022,li2023seg_survey}. For a single image defined on $\Omega \subset \mathbb{R}^2$, segmentation can be written as a mapping
\begin{equation}
S: \Omega \rightarrow \mathcal{L},
\end{equation}
where $\mathcal{L}$ is a set of labels. In \emph{semantic segmentation} (SS), $S(u,v)$ denotes the class label of pixel $(u,v)$. 
In \emph{instance segmentation} (IS), the output is a set of masks
$\{M_k\}_{k=1}^{K}$ with associated class labels $\{\ell_k\}$, where each $M_k:\Omega \rightarrow \{0,1\}$ corresponds to an object instance \cite{long2015fcn, he2017mask}.
 \emph{Panoptic segmentation} (PS) unifies SS and IS by predicting a single, non-overlapping labeling of pixels that distinguishes both ``stuff'' and countable object instances \cite{kirillov2019panoptic}.

These tasks extend naturally to videos. Given a sequence of
frames $\{I_t\}_{t=1}^{T}$, \emph{video semantic segmentation} (VSS) predicts per-frame semantic labels $\{S_t\}$, while \emph{video instance segmentation} (VIS) predicts instance masks with temporal consistency across frames, often requiring instance association (tracking) over time \cite{li2023seg_survey}. 
\emph{Video panoptic segmentation} (VPS) further combines VSS and VIS, producing temporally consistent panoptic outputs, and is commonly evaluated with metrics that account for both segmentation quality and tracking consistency~\cite{li2023seg_survey}. 
Closely related
settings include \emph{video object segmentation} (VOS)~\cite{perazzi2016benchmark}, which segments foreground objects across time under semi-supervised or unsupervised protocols, and \emph{referring} segmentation, which predicts masks grounded by natural language expressions in images (RIS)~\cite{wang2022cris} or videos (RVOS)~\cite{seo2020urvos}.
Segmentation maps are typically formed by dense prediction models that produce either hard labels or per-pixel class probabilities $P(u,v,\ell)$ (and, for instance-level tasks, a set of object masks and scores). 
A key advantage of segmentation is its explicit and interpretable structure: it converts low-level pixels into object- and region-level representations that support downstream reasoning and decision-making, and it can be fused naturally with geometry (\eg, depth, normals) and language in multimodal systems~\cite{li2023seg_survey}. 
Consequently, segmentation maps are central to applications such as autonomous driving,
medical image analysis, robotic manipulation, visual editing, and human-computer interaction.


\subsection{Depth Maps}
\label{subsec:p-2-sec-4-2D-depth-maps}
A depth map is a geometric representation that augments a 2D image with per-pixel distance information, describing the spatial layout of a scene from a given viewpoint~\cite{szeliski2022}. 
Formally, given an image domain $\Omega \subset \mathbb{R}^2$, a depth map is a scalar function
\begin{equation}
D: \Omega \rightarrow \mathbb{R},
\end{equation}
where $D(u,v)$ denotes the distance along the camera's optical axis to the first visible surface at pixel $(u,v)$; some sensors instead report radial range along the viewing ray.
Related representations store inverse depth or disparity, which must be converted to the appropriate metric depth or range before geometric back-projection~\cite{hartley2003multiple}.
Depth maps can be obtained through multiple mechanisms, including stereo vision via triangulation~\cite{hartley2003multiple}, active sensors such as structured-light or time-of-flight cameras~\cite{foix2011tof}, and learning-based monocular depth estimation~\cite{eigen2014depth}. 
When combined with known camera intrinsics, a depth map enables the back-projection of image pixels into 3D space, thereby serving as a bridge between 2D image representations and 3D geometry.
A key advantage of depth maps lies in their compact, image-aligned structure, which allows them to be processed using standard convolutional architectures and easily fused with the RGB appearance information. 
Compared to volumetric representations, depth maps are memory-efficient while still encoding meaningful geometric cues. 
As a result, they are widely used in applications such as 3D reconstruction, novel view synthesis~\cite{mildenhall2021nerf}, robotic navigation, and autonomous driving~\cite{KITTI-GeigerLU12}.


\subsection{Surface Normal Maps}
\label{subsec:p-2-sec-4-2D-surface-normal-maps}
A surface normal map is a geometric representation that encodes local surface orientation at each pixel of an image. Given an image domain
$\Omega \subset \mathbb{R}^2$, a surface normal map assigns to each pixel a unit vector
\begin{equation}
\mathbf{N}: \Omega \rightarrow \mathbb{S}^2,
\end{equation}
where $\mathbf{N}(u,v) = (n_x, n_y, n_z)$ represents the direction of the surface normal in camera or world coordinates, and $\mathbb{S}^2$ denotes the unit sphere. Surface normal maps can be derived from depth maps or 3D geometry via local differentiation, or directly estimated from RGB images using learning-based methods~\cite{barron2015shape}.
Classical approaches to normal estimation rely on photometric stereo or geometric reconstruction~\cite{woodham1980photometric}, while modern methods employ CNNs to regress surface normals from monocular images~\cite{eigen2015predicting}. 
Regardless of the estimation procedure, surface normal maps remain image-aligned and
encode rich local geometric information that complements depth and segmentation cues.
Surface normal maps can capture fine-grained shape and orientation details that may be ambiguous or noisy in raw depth estimates.
Because surface normals are invariant to absolute depth scale, they provide robust geometric constraints for scene understanding. 
Consequently, surface normal maps are widely used in applications such as 3D reconstruction, object pose estimation, scene layout analysis, realistic rendering, and as auxiliary supervision in multi-task learning~\cite{eigen2015predicting}.


\subsection{Edge Maps}
\label{subsec:p-2-sec-4-2D-edge-maps}
An edge map is a structural representation that highlights salient intensity or color discontinuities in an image, which often correspond to object boundaries and significant geometric transitions~\cite{canny1986}.
Given an image $I: \Omega \rightarrow \mathbb{R}^C$, an edge map is typically constructed by measuring local image gradients, for example,
\begin{equation}
E(u,v) = \|\nabla I(u,v)\|,
\end{equation}
where $\nabla I$ denotes the spatial gradient operator.
Classical edge detection methods rely on hand-crafted operators such as Sobel, Laplacian of Gaussian, and the Canny detector~\cite{canny1986}, while modern learning-based approaches directly predict the edge probabilities using convolutional neural networks~\cite{xie2015hed}.
Edge maps may be binary, indicating the presence or absence of an edge, or real-valued, encoding edge strength, orientation, or confidence.
The primary advantage of edge maps is their ability to significantly reduce visual complexity while preserving essential geometric structure.
This abstraction makes them relatively robust to changes in illumination and appearance. 
Consequently, edge maps are widely used in image
segmentation, contour detection, shape analysis, line drawing, and as intermediate cues for higher-level tasks such as object recognition and depth estimation~\cite{szeliski2022}.


\subsection{Optical Flow Maps}
\label{subsec:p-2-sec-4-2D-optical-flow-maps}
An optical flow map is a motion representation that encodes the apparent 2D displacement of pixels between consecutive frames in a video sequence~\cite{horn1981optical}. 
Given two image frames $I_t$ and $I_{t+1}$, the optical flow is defined as a vector field
\begin{equation}
\mathbf{F}(u,v) = (f_x(u,v), f_y(u,v)),
\end{equation}
where $(f_x, f_y)$ describes the horizontal and vertical displacement of pixel $(u,v)$ from time $t$ to $t+1$.
Classical optical flow estimation relies on assumptions such as
brightness constancy and spatial smoothness, leading to variational formulations~\cite{horn1981optical,brox2004high}. More recently, deep neural networks have been proposed to directly regress dense flow fields from image pairs or short video clips~\cite{dosovitskiy2015flownet,teed2020raft}.
Regardless of the estimation method, the resulting optical flow map remains image-aligned, preserving spatial correspondence with the input frames.
The main advantage of optical flow maps is their ability to capture rich dynamic information about scene motion, including object motion, camera motion, and motion parallax, while remaining computationally efficient.
As a result, optical flow is a core representation in video analysis and
robotic perception tasks such as motion segmentation, action
recognition, video interpolation, and visual odometry.



\subsection{2D Keypoints}
\label{subsec:p-2-sec-4-2D-2d-keypoints}
A 2D keypoint set is a sparse image representation that encodes the image-plane locations of a finite set of semantically or geometrically meaningful points, such as corners, joints, or landmarks~\cite{szeliski2022}.
Unlike dense maps that assign a value to every pixel, keypoint representations summarize visual content with a small number of salient points, trading pixel-level coverage for compactness and geometric explicitness.
Formally, given an image defined on a domain
$\Omega \subset \mathbb{R}^2$, a keypoint representation is a set of $K$ points
\begin{equation}
\mathcal{K} = \{(x_k, y_k)\}_{k=1}^{K}, \quad (x_k, y_k) \in \Omega,
\end{equation}
where each keypoint corresponds to a specific semantic or geometric entity. 
In learning-based formulations, keypoints are often represented implicitly as confidence or heat maps:
\begin{equation}
H_k : \Omega \rightarrow [0,1],
\end{equation}
with the keypoint location obtained by taking the argmax of $H_k$ or the expected image coordinates under its normalized weights \cite{newell2016stacked}. 
2D keypoints can be obtained using classical feature detectors, such as Harris corners, SIFT, and SURF \cite{harris1988combined,lowe2004distinctive}, or through modern deep neural networks that directly regress keypoint locations or predict dense heat maps \cite{toshev2014deeppose,newell2016stacked}.
Depending on the task, keypoints may represent generic interest points, object-specific landmarks, or articulated structures such as the human body joints.
The detection \emph{task} that produces such a set, including its extension to landmark detection and human pose estimation, is discussed in detail in Section~\ref{sec:keypoint_landmark} of Chapter~\ref{sec:perception}.
A key advantage of 2D keypoint representations is their sparsity and compactness: they significantly reduce the dimensionality of visual data while preserving critical structural information. 
This makes keypoints computationally efficient, robust to background clutter, and well-suited for geometric reasoning and correspondence matching.
Human pose estimation is a specialized and highly structured instance of 2D keypoint detection that focuses on localizing a predefined set of anatomical landmarks of the human body, such as joints and facial keypoints, in images or videos \cite{zheng2022pose}. 
In the 2D setting, the task aims to estimate the image coordinates of body joints, typically represented as a set of keypoints
$\mathcal{P} = \{(x_j, y_j)\}_{j=1}^{J}$, where $J$ denotes the number of joints in a human skeletal model. 
These keypoints are not independent: they are constrained by the underlying human body structure, which is often modeled using kinematic trees, graphical models, or learned
structural priors.






\section{3D}
\label{sec:rep_3d}

\FramedInitial{T}{he} 3D representations lift the 2D pixel grid into the ambient space the scene actually occupies, replacing image-plane appearance with explicit geometry and topology.
They are the substrate on which downstream tasks that require metric reasoning, occlusion handling, viewpoint manipulation, or physical simulation are built, from autonomous driving and robotic manipulation to medical visualization and immersive graphics.
Compared with the depth and surface normal maps of Section~\ref{sec:rep_2d}, which still inhabit the image grid, 3D representations expose geometry directly and can therefore describe parts of the scene that are not visible in any single view.

We organize this section around the eight 3D representations that dominate current practice (Figure~\ref{fig:p2_s4_2_3d_representation_1}, Table~\ref{tab:p2_s4_1_3d_comparison}): RGB-D images, voxels, point clouds, meshes, implicit fields, neural radiance fields (NeRF), 3D Gaussian Splatting (3DGS), and hybrids of the above.
They differ in what they store (surfaces vs.\ volumes, explicit primitives vs.\ neural functions), how they trade off fidelity against memory and rendering speed, and how amenable they are to editing and to gradient-based optimization.

\begin{itemize}[leftmargin=2em,itemsep=0.35em]
    \item \textbf{RGB-D Image}. Pairs a standard RGB image with a per-pixel depth map; view-dependent and cannot capture occluded surfaces from a single frame.
    \item \textbf{Voxel}. A uniform 3D grid of volumetric elements; conceptually simple and well suited to inherently volumetric data, but memory grows as $O(N^3)$ for $N$ cells per axis.
    \item \textbf{Point Cloud}. An unstructured set of 3D points, typically the native output of LiDAR or photogrammetry; faithful to the raw measurement but lacks surface connectivity.
    \item \textbf{Mesh}. Surfaces defined by vertices and faces; the dominant representation in graphics for its GPU-friendliness, at the cost of a fixed connectivity that complicates topology change.
    \item \textbf{Implicit Representation}. Defines surfaces as the level set of a continuous function (\eg, Signed Distance Functions); resolution-agnostic and topology-flexible, but slow to render.
    \item \textbf{NeRF (Neural Radiance Field)}. A neural function over a 5D input domain (3D location + 2D viewing direction) that supports state-of-the-art novel-view synthesis.
    \item \textbf{3DGS (3D Gaussian Splatting)}. Represents a scene as anisotropic 3D Gaussian primitives rendered through a differentiable splatting rasteriser, reaching NeRF-level quality at real-time speeds.
    \item \textbf{Hybrid Methods}. Combinations of the above (\eg, DMTet, Tri-plane) that pair explicit primitives with implicit decoders to balance fidelity, efficiency, and gradient-based optimizability.
\end{itemize}

\begin{figure}[!ht]
  \centering
  \includegraphics[width=0.99\linewidth]{figures/part2/3d_representation.pdf}
  \caption{Examples of 3D representations: RGB-D image, voxel grid, point cloud, triangle mesh, and a signed distance field (SDF) slice.}
  \label{fig:p2_s4_2_3d_representation_1}
\end{figure}


% \bigskip

\begin{table}[!t]
\centering
% \small
% \setlength{\tabcolsep}{4pt} % Adjust column spacing for a better fit
\caption{Qualitative comparison of 3D representations (relative ratings, not measured values).}
\label{tab:p2_s4_1_3d_comparison}
% \fontsize{9}{11}\selectfont
% \setlength{\tabcolsep}{2.5mm}
\resizebox{\linewidth}{!}{
\begin{tabular}{l|ccccc}
\toprule
\textbf{Representation} & \textbf{Fidelity} & \textbf{Memory Cost} & \textbf{Render Speed} & \textbf{Topology Flexibility} & \textbf{Editability} \\
\midrule
\textbf{RGB-D Image}    & Low (2.5D)  & Low                 & High                  & View-dependent           & Low          \\
\textbf{Voxel}          & Low--Medium & High ($O(N^3)$)     & Medium                & High                     & Medium       \\
\textbf{Point Cloud}    & Medium      & Medium (scalable)   & High                  & High                     & Medium       \\
\textbf{Mesh}           & High        & Low (GPU-friendly)  & Very High             & Low (fixed connectivity) & High         \\
\textbf{Implicit (SDF)} & High        & Low                 & Low (volume query)    & High                     & Low          \\
\textbf{NeRF}           & Very High   & Low                 & Low (ray marching)    & High                     & Very Low     \\
\textbf{3DGS}           & Very High   & Medium              & Very High (real-time) & High                     & Medium       \\
\textbf{Hybrid}         & Very High   & Medium              & High                  & Adaptive                 & Medium--High \\
\bottomrule
\end{tabular}
}
\end{table}


\subsection{RGB-D Images}
\label{sec:rep_3d_rgbd}

RGB-D images are a popular 2.5D representation that captures both photometric and geometric information from a single viewpoint.
This data format consists of two co-registered images: a standard RGB color image providing appearance and texture details, and a corresponding depth map where each pixel value encodes the camera-space depth of the corresponding scene point.
Formally, it can be defined as a function mapping a 2D image grid $\Omega$ to a 4D space combining color and depth:
\begin{equation}
I_{rgbd}: \Omega \rightarrow \mathbb{R}^3 \times \mathbb{R}_{\geq 0},
\end{equation}
where the first three channels represent RGB color and the fourth channel encodes the depth of the corresponding scene point. In practice, RGB-D data is stored as a tensor $I \in \mathbb{R}^{H \times W \times 4}$.

The primary advantage of RGB-D data lies in its efficiency and accessibility.
It is captured by low-cost sensors such as Microsoft Kinect or Intel RealSense, and its image-like grid structure enables the direct application of powerful 2D convolutional neural networks (CNNs) to tasks such as semantic segmentation and object detection~\cite{dai2016r,couprie2013indoor,vu2019dada}.
This fusion of color and depth provides a rich, multi-modal input for scene understanding.
However, a significant limitation is its view-dependent nature; an RGB-D image represents only a partial view of the 3D world, suffering from occlusions and providing no information about surfaces hidden from the camera.
Despite this, its real-time acquisition and processing capabilities make it indispensable in robotics for navigation~\cite{desouza2002vision} and grasping~\cite{bicchi2000robotic,saxena2006robotic,zhao2025towards}, in augmented reality (AR) for scene interaction, and as a foundational input for 3D reconstruction pipelines that fuse multiple frames into a complete model~\cite{newcombe2011kinectfusion}.







\subsection{Voxel Grids}
\label{sec:rep_3d_voxel}

The voxel grid is a foundational 3D representation that discretizes space into a uniform grid of volumetric elements, or \textit{voxels} (a portmanteau of ``volume'' and ``pixel'').
Analogous to a pixel in 2D, a voxel grid can be modeled as a function mapping discrete 3D coordinates to a set of attributes:
\begin{equation}
G: \{1, \dots, D\} \times \{1, \dots, H\} \times \{1, \dots, W\} \rightarrow \mathcal{A},
\end{equation}
where $D, H, W$ are the grid dimensions and $\mathcal{A}$ can represent occupancy, color, density, or material properties, thereby enabling a detailed volumetric description of an object.

A key advantage of this representation is its ability to model complex geometries and internal structures, which are often challenging for surface-based methods like polygon meshes to capture.
The conceptual simplicity of voxels, which can be seen as the three-dimensional extension of 2D pixels, facilitates straightforward implementation and manipulation.
Consequently, voxel grids are natively suited for applications involving inherently volumetric data, such as medical imaging (\eg, MRI and CT scans)~\cite{cciccek20163d}, where biological tissues are naturally captured as a 3D array of values.
They are also widely adopted in computer graphics and gaming for tasks like procedural generation and creating destructible environments, as popularized by games like \textit{Minecraft}.

\subsection{Point Clouds}
\label{sec:rep_3d_point}

Point clouds offer a direct and unstructured 3D representation, defined as a set of points in 3D space:
\begin{equation}
P = \{p_i\}_{i=1}^N, \quad p_i = (x_i, y_i, z_i, \mathbf{a}_i) \in \mathbb{R}^3 \times \mathbb{R}^C.
\end{equation}
Each point $p_i$ consists of its geometric coordinates $(x_i, y_i, z_i)$ and an optional $C$-dimensional attribute vector $\mathbf{a}_i$ encoding properties such as color, intensity, or surface normals. They are the native output of active 3D scanning technologies like LiDAR (Light Detection and Ranging) and structured-light scanners~\cite{behroozpour2017lidar}, as well as passive methods like photogrammetry~\cite{schonberger2016structure}.

The primary strength of point clouds lies in their raw and unstructured nature; they represent geometry without imposing any explicit topological connectivity (unlike meshes) or spatial discretization (unlike voxels).
This makes them a faithful and often highly precise representation of the captured data.
However, this very lack of inherent structure poses significant challenges for downstream processing.
The unordered and potentially sparse distribution of points makes tasks like feature extraction, segmentation, and reconstruction non-trivial, necessitating specialized algorithms and deep learning architectures (\eg, PointNet~\cite{qi2017pointnet} and its successors~\cite{qi2017pointnetplus}) that are designed to operate directly on irregular point sets.


\subsection{Meshes}
\label{sec:rep_3d_mesh}

A mesh is a canonical surface representation in computer graphics, defining an object's boundary with a set of geometric and topological primitives. Formally, a triangle mesh is a tuple $M = (V, F)$, where
\begin{equation}
V = \{v_i\}_{i=1}^{|V|} \subset \mathbb{R}^3, \; F\subset V \times V \times V.
\end{equation}
Here, $V$ is a set of vertices representing points in 3D space, and $F$ is a set of faces (typically triangles) that define connectivity between vertices.

The key characteristic of a mesh is its explicit topological structure, which encodes the connectivity between these elements. This structure strikes an effective balance between descriptive power and computational efficiency. Its prevalence is due to several key advantages: meshes can represent complex, smooth surfaces with arbitrary fidelity through adaptive tessellation (\ie, varying the polygon density); moreover, this representation is natively supported by modern Graphics Processing Units (GPUs), enabling highly efficient real-time rendering. The explicit connectivity also facilitates the straightforward computation of essential surface properties like normals and curvature, which are critical for realistic lighting and physical analysis. However, as a purely surface-based representation, standard meshes are ill-suited for modeling volumetric or internal structures, and representing objects with complex or changing topologies can be algorithmically challenging.








\subsection{Implicit Representations}
\label{sec:rep_3d_implicit}

Implicit representations define 3D geometry not as a discrete set of primitives (\eg, points or polygons), but as the level set of a continuous function:
\begin{equation}
f: \mathbb{R}^3 \to \mathbb{R}.
\end{equation}
The surface $\mathcal{S}$ is implicitly defined as the set of points where the function evaluates to a specific value, typically zero:
\begin{equation}
\mathcal{S} = \{\mathbf{x} \in \mathbb{R}^3 | f(\mathbf{x}) = 0\}.
\end{equation}

A particularly powerful and popular form of this is the \textbf{Signed Distance Function (SDF)}.
An SDF maps any spatial coordinate $\mathbf{x}$ to the closest distance to the surface, where the sign indicates whether the point is inside ($f(\mathbf{x}) < 0$) or outside ($f(\mathbf{x}) > 0$) the object.
This formulation inherently guarantees watertight, closed surfaces and provides a smooth, continuous description of the geometry, making it resolution-agnostic.
A key advantage of implicit fields is their ability to elegantly handle complex and dynamic topologies; operations like merging, splitting, or creating holes do not require complex and error-prone re-meshing procedures.
This makes them ideal for physical simulations, such as fluid dynamics and shape manipulation.

This paradigm has experienced a major resurgence with the rise of \textbf{neural fields} or \textbf{neural implicit representations}. Instead of storing the function in a grid, a neural network is trained to approximate it~\cite{mescheder2019occupancy} (\eg, DeepSDF~\cite{park2019deepsdf}).
This approach is memory-efficient and can learn continuous, high-fidelity shape priors from data.
Extensions such as Neural Radiance Fields (NeRF) further generalize implicit fields to jointly model geometry and appearance, and are discussed in detail in Section~\ref{sec:rep_3d_nerf}.
To visualize or use implicit surfaces in traditional rendering pipelines, an explicit mesh representation is often extracted via algorithms like Marching Cubes~\cite{lorensen1998marching}.






\begin{figure}[!ht]
  \centering
  \includegraphics[width=\linewidth]{figures/part2/3D_representations_2.pdf}
  % \vspace{-7pt}
  \caption{Neural and hybrid 3D representations: (a) NeRF, (b) 3DGS, and (c) hybrid methods (DMTet and Tri-plane).}
  \label{fig:p2_s4_2_3d_representation_2}
\end{figure}


\subsection{Neural Radiance Fields}
\label{sec:rep_3d_nerf}

Neural Radiance Fields (NeRF) represent a paradigm shift in novel view synthesis and implicit 3D scene representation \cite{mildenhall2021nerf}.
The core innovation is to represent a static scene as a continuous function with a 5D input domain (3D location $+$ 2D viewing direction), parameterized by a simple Multi-Layer Perceptron (MLP).
This neural network, denoted as $F_\Theta$, maps a 3D spatial coordinate $\mathbf{x}=(x,y,z)$ and a 2D viewing direction $\mathbf{d}=(\theta, \phi)$ to a volume density $\sigma$ and an emitted color $\mathbf{c}=(r,g,b)$.
A crucial design choice is that density $\sigma$ depends only on the location $\mathbf{x}$, while the color $\mathbf{c}$ depends on both location and viewing direction $(\mathbf{x}, \mathbf{d})$.
This allows NeRF to capture complex, non-Lambertian surface effects such as specular reflections and translucency.

To render an image from a novel viewpoint, rays are cast through the scene, and their colors are computed by integrating the predicted color and density values along each ray using principles of classical volume rendering and numerical quadrature.
By optimizing the network's weights to minimize the photometric reconstruction error between rendered and ground-truth images, NeRF can learn a highly detailed and coherent representation of the scene from only a sparse set of posed input images.
The resulting state-of-the-art photorealism has established NeRF as a foundational technique, inspiring a vast body of follow-up research aimed at improving its speed, scalability, and applicability to dynamic scenes~\cite{muller2022instant,barron2021mip,barron2023zip,multinerf2022}.

\subsection{3D Gaussian Splatting}
\label{sec:rep_3d_gaussian}

3D Gaussian Splatting (3DGS) has recently emerged as a groundbreaking alternative to neural radiance fields, achieving state-of-the-art quality for novel view synthesis while enabling real-time rendering \cite{kerbl20233d}.
Unlike the implicit representation of NeRF, 3DGS represents a scene as a set of $N$ anisotropic 3D Gaussian primitives, each parameterized by:
\begin{itemize}[leftmargin=2em,itemsep=0.35em]
    \item \textbf{Position} $\boldsymbol{\mu} \in \mathbb{R}^3$: the 3D center of the Gaussian;
    \item \textbf{Covariance} $\boldsymbol{\Sigma} \in \mathbb{R}^{3\times 3}$: shape and orientation, allowing thin planes through to volumetric wisps;
    \item \textbf{Color}: Spherical Harmonics (SH) coefficients capturing view-dependent appearance;
    \item \textbf{Opacity} $\alpha \in [0,1]$: opacity ($0$ is transparent and $1$ is opaque).
\end{itemize}
The density contributed by a single primitive at a query point $\mathbf{x} \in \mathbb{R}^3$ is
\begin{equation}
G(\mathbf{x}; \boldsymbol{\mu}, \boldsymbol{\Sigma}) = \exp\!\left(-\tfrac{1}{2} (\mathbf{x}-\boldsymbol{\mu})^{\top} \boldsymbol{\Sigma}^{-1} (\mathbf{x}-\boldsymbol{\mu}) \right).
\end{equation}
$\boldsymbol{\Sigma}$ is decomposed into a scaling matrix $\mathcal{S}$ and rotation matrix $\mathcal{R}$, such that
\begin{equation}
\boldsymbol{\Sigma} = \mathcal{R}\mathcal{S}\mathcal{S}^T\mathcal{R}^T,
\end{equation}
with per-primitive scaling and rotation stored as $s \in \mathbb{R}^{N \times 3}$ and quaternions $r \in \mathbb{R}^{N \times 4}$. The key to 3DGS's real-time performance lies in the differentiable splatting rasteriser introduced together with the representation~\cite{kerbl20233d}, which extends earlier differentiable rasterization work on meshes~\cite{liu2019soft}.
Instead of costly per-pixel ray marching, 3DGS projects (or ``splats'') these 3D Gaussians onto the 2D image plane, creating 2D Gaussian splats.
These splats are then sorted and alpha-blended in a highly-parallelized, feed-forward process that is exceptionally well-suited for modern GPU architectures.
Differentiable splatting applies a viewing transform $W$ and Jacobian $J$ to compute the transformed covariance $\boldsymbol{\Sigma}^{\prime} = JW\boldsymbol{\Sigma} W^TJ^T$, enabling novel view rendering. Each pixel color $\mathcal{C}$ is obtained by blending the $N$ overlapping primitives along the ray:
\begin{equation}
\mathcal{C} = \sum_{i=1}^{N}c_i \alpha_i \prod_{j=1}^{i-1} (1-\alpha_j),
\end{equation}
where $c_i$ and $\alpha_i$ denote the color and opacity of the $i$-th primitive.
The entire model, initialized from a sparse point cloud (typically from Structure-from-Motion), is optimized end-to-end against training images.
The optimization process also includes an adaptive density control mechanism to dynamically add, split, or prune Gaussians to best represent the scene's geometry.
This combination of an explicit, flexible primitive with a fast, differentiable rasterizer has set a new standard for performance and quality in 3D scene reconstruction.





\subsection{Hybrid Methods}
\label{sec:rep_3d_hybrid}
To harness the complementary strengths of different 3D representations, hybrid methods have emerged as a powerful paradigm.
These approaches combine two or more representations, integrating the high-fidelity and continuous nature of implicit functions with the computational efficiency and explicit structure of primitives like meshes or voxels.
This section reviews two influential hybrid techniques: Deep Marching Tetrahedra (DMTet) and Tri-plane.

\paragraph{Deep Marching Tetrahedra (DMTet).}
Deep Marching Tetrahedra (DMTet)~\cite{shen2021deep} is an innovative hybrid technique that bridges learned implicit fields with explicit polygonal meshes. 
Its core concept involves using a neural network to predict an implicit representation of a shape, such as a Signed Distance Function (SDF)~\cite{vicini2022differentiable} or an occupancy field. 
A differentiable version of the marching tetrahedra algorithm subsequently extracts a high-quality, explicit triangle mesh surface from this implicit representation.

The primary advantage of DMTet lies in its end-to-end differentiability. 
The entire pipeline, from the implicit field prediction to the final mesh generation, allows gradients to flow from a 2D-image-based loss function back to the 3D shape parameters.
This enables direct optimization of the mesh geometry using 2D supervision, leading to highly detailed and accurate models.
Furthermore, DMTet can learn to produce an adaptive mesh, allocating higher tetrahedral density to regions of fine detail, thus optimizing computational resources while maintaining high fidelity.


\paragraph{Tri-plane.}
The Tri-plane representation~\cite{chan2022efficient} offers an efficient, factorized approach to modeling 3D scenes and objects. Instead of encoding information in a dense 3D volume, it decomposes the 3D space into three orthogonal 2D feature planes (typically aligned with the XY, YZ, and XZ axes). This design significantly reduces the memory and computational complexity associated with 3D representations.

The working principle is to derive the feature for any 3D point by projecting it onto the three 2D planes, querying the features at the corresponding projected coordinates, and aggregating them (\eg, through summation or concatenation). 
This aggregated feature vector is then fed into a small MLP to decode the final property of the point, such as its color, density, or SDF value.
This ``explicit-grid-plus-implicit-decoder'' structure is the essence of its hybrid nature.
It benefits from the efficiency of 2D grid-based feature storage, which can be processed by powerful 2D CNNs, while avoiding the prohibitive memory cost of a full 3D grid.


\begin{table}[!t]
\centering
\caption{Representative methods for each 3D representation family, chosen to show its historical progression and dominant design choices.}
\label{tab:p2_s4_3d_representative_methods}
% \small
\fontsize{9}{11}\selectfont
\setlength{\tabcolsep}{4pt}
\begin{tabular}{L{0.19\textwidth}L{0.4\textwidth} L{0.4\textwidth}}
\toprule
\textbf{Representation} & \textbf{Representative Methods} & \textbf{Key Idea / Contribution} \\
\midrule
\textbf{RGB-D Image} & KinectFusion~\cite{newcombe2011kinectfusion}, DADA~\cite{vu2019dada}, SemanticFusion~\cite{mccormac2017semanticfusion} & Fuses color and per-pixel depth for real-time reconstruction and semantic understanding. \\
\textbf{Voxel Grid} & VoxNet~\cite{maturana2015voxnet}, 3D U-Net~\cite{cciccek20163d}, OGN~\cite{tatarchenko2017octree} & Direct 3D CNNs on regular grids; octree variants reduce memory for high resolution. \\
\textbf{Point Cloud} & PointNet~\cite{qi2017pointnet}, PointNet++~\cite{qi2017pointnetplus}, Point Transformer~\cite{zhao2021point}, PointNeXt~\cite{qian2022pointnext} & Permutation-invariant and attention-based architectures for unordered point sets. \\
\textbf{Mesh} & MeshCNN~\cite{hanocka2019meshcnn}, Pixel2Mesh~\cite{wang2018pixel2mesh}, AtlasNet~\cite{groueix2018papier} & Learning on mesh edges and single-image mesh generation from data. \\
\textbf{Implicit (SDF/occupancy)} & Occupancy Networks~\cite{mescheder2019occupancy}, DeepSDF~\cite{park2019deepsdf}, IM-Net~\cite{chen2019imnet}, NeuS~\cite{wang2021neus} & Neural networks encode continuous surfaces via occupancy or signed-distance fields. \\
\textbf{NeRF} & NeRF~\cite{mildenhall2021nerf}, Plenoxels~\cite{fridovich2022plenoxels}, Instant-NGP~\cite{muller2022instant}, Mip-NeRF 360~\cite{barron2022mip} & 5D radiance fields; later works accelerate training and reduce aliasing. \\
\textbf{3DGS} & 3D Gaussian Splatting~\cite{kerbl20233d}, Mip-Splatting~\cite{yu2024mip}, Scaffold-GS~\cite{lu2024scaffold}, 4DGS~\cite{wu20244d} & Anisotropic Gaussian primitives with differentiable rasterization for real-time rendering. \\
\textbf{Hybrid} & DMTet~\cite{shen2021deep}, Tri-plane~\cite{chan2022efficient}, NeRFactor~\cite{zhang2021nerfactor} & Combines explicit grids or meshes with compact implicit decoders. \\
\bottomrule
\end{tabular}
\end{table}

% \subsection{Comparison, Applications, and Challenges}
% \label{sec:rep_3d_summary}

% Table~\ref{tab:p2_s4_1_3d_comparison} highlights the complementary trade-offs among the eight representations surveyed above, while Table~\ref{tab:p2_s4_3d_representative_methods} collects representative methods.
% RGB-D images, voxels, point clouds, and meshes remain the workhorses of applications that prize acquisition simplicity, volumetric modeling, raw geometric fidelity, or GPU rendering, respectively.
% Neural representations, NeRF and 3DGS in particular, have shifted the frontier toward photorealistic novel-view synthesis from sparse images, but they do so at the cost of editability and, in the case of NeRF, rendering speed.
% Hybrid methods attempt to occupy the middle ground by pairing explicit primitives with compact neural decoders.

% These trade-offs manifest differently across application domains.
% In \textbf{robotics}, RGB-D sensors and point clouds dominate perception and grasping pipelines; meshes and implicit fields are common for collision checking and simulation.
% In \textbf{autonomous driving}, LiDAR point clouds are the canonical input, while voxel grids and bird's-eye-view feature maps serve as intermediate representations for detection and prediction.
% In \textbf{graphics and virtual production}, meshes remain the standard for real-time rendering, while NeRF and 3DGS are increasingly used for view synthesis, digital humans, and scene relighting.
% In \textbf{medical imaging}, voxel grids are natural for CT and MRI data, and implicit surfaces are used for organ and tumor segmentation.

% Several open challenges cut across all of these representations.
% The most pervasive is \textbf{scalability}: as scenes grow from single objects to rooms, buildings, and outdoor environments, memory and computation become bottlenecks for voxels, point clouds, and dense Gaussian primitives alike.
% Even where memory is not the limiting factor, \textbf{editability and composability} lag behind: neural fields and 3DGS are difficult to edit in the intuitive way meshes or voxels allow, and combining several reconstructed assets into a coherent scene remains nontrivial.
% Time compounds both problems, since most representations surveyed here assume static geometry, and extending them to \textbf{dynamic and deformable scenes} (articulated objects, fluids, general deformable media) is still an active research direction.
% Underlying all of these is \textbf{cross-representation conversion} (extracting meshes from neural fields, compressing point clouds into Gaussians, or turning implicit surfaces into editable primitives), which remains a practical barrier to deploying these representations in downstream pipelines.
% Progress will therefore depend not only on better individual representations, but on principled ways to compose and convert between them.



\section{4D}
\label{sec:rep_4d}

\FramedInitial{T}{he} 4D representations extend the 3D representations of Section~\ref{sec:rep_3d} into the temporal dimension, modeling not only the geometry and appearance of a scene at a single instant but also how both evolve over time.
They are the substrate for tasks that involve dynamic content (moving objects, deformable bodies, fluid simulation, novel-view synthesis of casually captured videos), where a sequence of independent 3D reconstructions would either fail to maintain temporal coherence or be prohibitively expensive to recompute frame by frame.

We organize this section around five 4D representations (Figure~\ref{fig:p2_s4_3_4d_representations}), each obtained by adding time to one of the 3D representations of Section~\ref{sec:rep_3d}: RGBD video, point clouds, meshes, neural radiance fields, and deformable 3D Gaussian Splatting.
The voxel and implicit-field families of Section~\ref{sec:rep_3d} admit analogous time extensions but are less commonly used in the 4D literature, and we therefore omit them as standalone subsections here.



\begin{figure}[!th]
    \centering
    \includegraphics[width=\linewidth]{figures/part2/4D_representations.pdf}
     % \vspace{-7pt}
    \caption{Five 4D representations, each a 3D primitive evolving over time: (a) RGBD video, (b) 4D point cloud, (c) deformable mesh, (d) dynamic NeRF, and (e) deformable 3DGS. The RGBD example uses changing camera poses across time to observe the scene from different viewpoints; the top-view camera icons use the same azimuths as the paired RGB and depth renderings. The point-cloud panel distinguishes independent frame-wise samples from persistent point trajectories. Mesh colors indicate vertex displacement; re-extraction is an alternative to preserving vertex and face identities. In (d), rows are views and columns are times. In (e), $R$ and $S$ abbreviate the time-dependent $R_i(t)$ and $S_i(t)$; constant opacity and color describe the illustrated parameterization. All scenes are controlled synthetic illustrations, not experimental reconstructions.}
    \label{fig:p2_s4_3_4d_representations}
\end{figure}



\subsection{RGBD Video}
\label{sec:rep_4d_rgbd}
RGBD videos $\mathcal{V}=\{(I_t, D_t)\}_{t=1}^T$, with $I_t \in \mathbb{R}^{H \times W \times 3}$ and $D_t \in \mathbb{R}^{H \times W}$, are a straightforward extension of RGBD images (see Section~\ref{sec:rep_3d_rgbd}) to the temporal domain.
Equivalently, they can be stored as a 4D tensor $\mathcal{V} \in \mathbb{R}^{T \times H \times W \times 4}$, with the last channel holding depth.
RGBD videos can be captured using depth sensors, such as Microsoft Kinect or Intel RealSense, which provide synchronized RGB and depth data.
These videos are commonly used in applications such as motion capture, gesture recognition, and augmented reality.
There are also many learning-based methods~\cite{wang2023neural,hu2025depthcrafter,chen2025video} that directly predict consistent RGBD videos from monocular video sequences
$\mathcal{I}=\{I_t\}_{t=1}^T$.
However, such a representation may struggle to handle complex dynamic scenes with significant occlusions or topology changes.
If we want to further reconstruct a unified 4D scene in a canonical space, we may also need to provide or estimate camera poses
$
\mathcal{P}=\{P_t\}_{t=1}^T
$
for each frame, as well as point trajectories over time to establish temporal correspondences, \eg via joint pose-trajectory optimization in recent video-based SfM systems~\cite{li2025megasam}.






\subsection{Point Cloud}
\label{sec:rep_4d_point}
Point clouds are another popular representation for 3D scenes (see Section~\ref{sec:rep_3d_point}), and they can be extended to 4D by incorporating temporal information.
A 4D point cloud can be represented as $\mathcal{P}=\{(\mathbf{p}_i(t), \mathbf{c}_i(t))\}_{i=1}^N$, where $\mathbf{p}_i(t) \in \mathbb{R}^3$
is the 3D position of point $i$ at time step $t$, and $\mathbf{c}_i(t) \in \mathbb{R}^3$ is the corresponding color.
To map an uncalibrated, continuous video sequence $\mathcal{I}=\{I_t\}_{t=1}^T, I_t\in \mathbb{R}^{H \times W \times 3}$, to a 4D point cloud representation, there are generally two types of representations:
(1) a sequence of point clouds $\mathcal{P}=\{\mathcal{P}_t\}_{t=1}^T$ in a fixed coordinate system, where each $\mathcal{P}_t=\{(\mathbf{p}_i(t), \mathbf{c}_i(t))\}_{i=1}^{H \times W}$ is an independent point cloud at time step $t$;
(2) a unified point cloud $\mathcal{P}=\{(\mathbf{p}_i(t), \mathbf{c}_i(t))\}_{i=1}^{N}$, where $N$ is the number of persistent point identities, and each point $i$ has a trajectory over time.
The first representation is easier to obtain, as it can be directly derived from the depth maps estimated from each frame using monocular depth estimation methods~\cite{li2025megasam}.
However, it may suffer from temporal inconsistencies and a lack of point correspondences across frames.
The second representation, in contrast, requires explicit temporal correspondences (either ground-truth tracks, optical-flow lifting~\cite{teed2020raft}, or jointly optimized point trajectories~\cite{li2025megasam}), and is therefore harder to acquire, but it offers persistent point identities and consistent trajectories that are useful for downstream 4D reconstruction, motion analysis, and dynamic NeRF/3DGS initialization.

\subsection{Mesh}
\label{sec:rep_4d_mesh}
We can extend the mesh representation (see Section~\ref{sec:rep_3d_mesh}) to 4D by representing a sequence of 3D meshes over time.
A 4D mesh is a sequence $\{\mathcal{M}_t\}_{t=1}^T$ with $\mathcal{M}_t = (\mathbf{V}_t, \mathbf{F}_t)$, where $\mathbf{V}_t \in \mathbb{R}^{N \times 3}$ and $\mathbf{F}_t \in \mathbb{N}^{M \times 3}$ are the vertices and faces of the mesh at time $t$, respectively.
To reconstruct a 4D mesh using deep learning, researchers typically sample a set of points on the mesh surface at the original time step $t=0$, and then learn a deformation field to model the motion of these points over time~\cite{cao2024motion2vecsets}.
Then, the deformed points at each time step can be used to reconstruct watertight manifold meshes using
Marching Cubes~\cite{lorensen1998marching} or Poisson Surface Reconstruction~\cite{kazhdan2006poisson}.








\subsection{NeRF}
\label{sec:rep_4d_nerf}
As discussed in Section~\ref{sec:rep_3d_nerf}, NeRF~\cite{mildenhall2021nerf} is a continuous, volume-based 3D representation that maps a 3D position $\mathbf{x} \in \mathbb{R}^3$ and a 2D viewing direction $\mathbf{d} \in \mathbb{S}^2$ to a color $\mathbf{c} \in \mathbb{R}^3$ and a volume density $\sigma \in \mathbb{R}_{\geq 0}$ via a learned function $F_\Theta:(\mathbf{x}, \mathbf{d}) \rightarrow (\mathbf{c}, \sigma)$.
This function is typically implemented using a neural network, which is trained on a set of input images and their corresponding camera poses.
To extend NeRF to dynamic scenes, we can introduce a temporal dimension $t$ to the input of the function, resulting in a 4D representation
$F_\Theta: (\mathbf{x}, \mathbf{d}, t) \rightarrow (\mathbf{c}_t, \sigma_t)$.
This allows the model to capture changes in the scene over time, enabling the rendering of novel views at different time steps.
In practice, the temporal dimension can be represented in various ways.
For example, a deformation network can take an observed position $\mathbf{x}$ and time $t$ as input and predict a displacement $\Delta\mathbf{x}(\mathbf{x},t)$ that maps it to a canonical position $\mathbf{x}_c=\mathbf{x}+\Delta\mathbf{x}(\mathbf{x},t)$~\cite{pumarola2021d}.
Alternatively, the deformation can be parameterized using a dense
$\text{SE}(3)$ field~\cite{park2021nerfies}, which models the rigid transformations (rotation and translation) of points in the scene over time.
Other approaches~\cite{gao2021dynamic} may also incorporate additional latent codes to represent different time steps, enabling the model to capture more complex temporal variations.







\subsection{Deformable 3D Gaussian Splatting}
\label{sec:rep_4d_gaussian}
Given a static scene representation $\mathcal{G}=\{(\alpha_i, \boldsymbol{\mu}_i, \boldsymbol{\Sigma}_i, c_i)\}_{i=1}^G$, which contains $G$ 3D Gaussians~\cite{kerbl20233d} (see Section~\ref{sec:rep_3d_gaussian}),
where $\alpha_i \in [0, 1]$ is the opacity, $\boldsymbol{\mu}_i\in\mathbb{R}^3$ is the mean (3D position), $\boldsymbol{\Sigma}_i\in\mathbb{R}^{3\times 3}$ is the covariance matrix, and $c_i: \mathbb{S}^2 \rightarrow \mathbb{R}^3$ is the color function that maps a viewing direction to an RGB color, we extend it to a \emph{deformable} 3D Gaussian Splatting representation for dynamic scenes by letting both the mean and the covariance vary in time: $\mathcal{G}(t)=\{(\alpha_i, \boldsymbol{\mu}_i(t), \boldsymbol{\Sigma}_i(t), c_i)\}_{i=1}^G$~\cite{wu20244d}.
The mean is parameterized by a deformation field, $\boldsymbol{\mu}_i(t)=\boldsymbol{\mu}_i(0)+\Delta\boldsymbol{\mu}_i(t)$, where $\boldsymbol{\mu}_i(0)$ is the initial position at $t=0$ and $\Delta\boldsymbol{\mu}_i(t)$ is the time-dependent offset.
The covariance is decomposed in the same way as in the static case, $\boldsymbol{\Sigma}_i(t)=\mathcal{R}_i(t)\mathcal{S}_i(t)\mathcal{S}_i(t)^{T}\mathcal{R}_i(t)^{T}$, with the rotation $\mathcal{R}_i(t)$ and scaling $\mathcal{S}_i(t)$ themselves outputs of the deformation field at time $t$.
The opacity $\alpha_i$ and color function $c_i$ are typically kept constant over time.
In practice, this means that the appearance of the scene does not change significantly over time, while the latest work~\cite{zhang2025gaussian} also explores time-varying colors and opacities.
This representation allows for efficient rendering of dynamic scenes by accounting for the temporal changes in the positions and shapes of the 3D Gaussians.




























\section{Language}
\label{sec:rep_language}

\FramedInitial{U}{nlike} the visual and geometric representations of the previous sections, which encode space explicitly through pixels, points, or surfaces, language represents space \emph{symbolically}, naming entities and stating the relations among them (\eg, ``the chair is next to the table'' or ``the store is across the street from the park'') and thereby capturing relational, topological, and functional structure that metric representations leave implicit. 
The question this section pursues is not how language \emph{describes} space, which is self-evident, but whether models trained only on language come to \emph{encode} it. 
Answering that requires first understanding how language representations are learned at all, so the section moves from the evolution of language representation learning, from masked-language pretraining to large autoregressive models, to the more surprising finding that such sequence-only training induces latent, spatially structured world models rather than mere text statistics. 
The scope here is language in isolation; how language is \emph{aligned} with visual, 3D, and audio modalities is deferred to Chapter~\ref{chapter:cross_modality}.



\begin{figure}[!ht]
    \centering
    \includegraphics[width=\linewidth]{figures/part2/language_representation.pdf}
    \caption{Spatial information implicit in language-only representations. \textit{Left}: natural-language descriptions of spatial relations. \textit{Right}: a probe on Othello-GPT hidden activations recovers the board state from move sequences alone.}
    \label{fig:p2_s4_1_language-representation}
\end{figure}


\subsection{The Evolution of Language Representation Learning}
\label{subsec:p-2-sec-4-the-evolution-of-language-representation-learning}

Language representation learning has evolved rapidly over the past decade, through masked-language pretraining (BERT~\cite{devlin2019bert} and RoBERTa~\cite{liu2019roberta}), permutation-based autoregressive pretraining (XLNet~\cite{yang2019xlnet}), denoising sequence-to-sequence pretraining (T5~\cite{raffel2020exploring}), and autoregressive large language models trained purely on next-token prediction, exemplified by GPT-2 and GPT-3~\cite{radford2019language,brown2020language}.
Across this progression, a recurring observation is that representation quality scales with model size, data size, and the sophistication of the pretraining objective~\cite{srivastava2022beyond}, and that pretrained representations encode richer structure than the surface text statistics they were trained on~\cite{tenney2019bert,hewitt2019structural}.
This last property, that language-only training produces representations carrying structure beyond text, is what makes language representations load-bearing for spatial intelligence, as the next subsection makes explicit.

\subsection{Language Representation Learning Induces Spatial Intelligence}
\label{subsec:p-2-sec-4-language-representation-learning-induces-spatial-intelligence}

Although LLMs are optimized for sequence prediction, recent work suggests that they can develop structured internal representations that resemble world models~\cite{li2021implicit,patel2022mapping,abdou2021can}. 
One particularly pertinent demonstration of this on spatial intelligence is Othello-GPT, a transformer trained solely on sequences of legal moves from the board game Othello~\cite{li2023emergent}.

Despite receiving no prior knowledge of the board layout or game rules, Othello-GPT achieves near-perfect accuracy in predicting legal moves (actions in a physical space). 
Probing experiments reveal that nonlinear~\cite{alain2016understanding,tenney2019bert,belinkov2016probing} classifiers can reliably reconstruct the entire board state from hidden activations~\cite{li2023emergent}. 
Crucially, interventional studies show that manipulating internal representations of specific tiles alters the model’s predictions in a way consistent with counterfactual board states, providing causal evidence of spatial encoding \cite{li2023emergent}.

These findings align with broader research indicating that LMs acquire implicit representations of abstract or grounded concepts. 
For example, models have been shown to encode syntactic structure~\cite{hewitt2019structural}, color spaces \cite{abdou2021can}, and even logical state-tracking~\cite{li2021implicit}. 
Work on chess transcripts similarly shows that LMs can track board positions without explicit supervision~\cite{toshniwal2021learning}. 
Compared to systems like AlphaZero, which are built with strong inductive biases and explicit access to board states~\cite{silver2018general,mcgrath2021acquisition}, Othello-GPT highlights the surprising extent to which sequence-only training can lead to emergent spatial reasoning.

In sum, it has been demonstrated that spatially structured representations can emerge in LMs trained only on symbolic sequences. 
This suggests that autoregressive sequence prediction is sufficient, in certain domains, to induce latent world models that encode spatial relations and can causally drive action predictions.








\section{Audio}
\label{sec:rep_audio}



\FramedInitial{A}{udio} is fundamentally tied to spatial information. 
Unlike vision, which directly captures the geometric layout of a scene, acoustic signals encode spatial structure only \emph{implicitly}, through the physical transformations imposed on sound during propagation.
Recovering space from sound is therefore a matter of reading these transformations backward: from the perceptual cues that human hearing exploits, to the computational representations that formalize them, to the learned models and cross-modal couplings that ground sound in the surrounding geometry. 
This section follows that arc.
Human auditory perception interprets the resulting acoustic cues, shaped by these transformations, to infer the spatial layout of the environment.
The primary cues are binaural: the \emph{Interaural Time Difference (ITD)} and the \emph{Interaural Level Difference (ILD)}, both arising from the geometry of the head and the spatial separation of the ears~\cite{macpherson2002listener,middlebrooks2015sound}.
ITD reflects the difference in arrival time between the two ears and is highly reliable for localizing low-frequency sounds whose long wavelengths exceed head size.
Conversely, ILD captures sound pressure level differences caused by head-shadowing effects and serves as a dominant cue for high-frequency localization where such attenuation is more pronounced.
However, binaural cues alone cannot fully resolve depth, elevation, or front-back ambiguities.

To address these limitations, the auditory system further exploits \emph{monaural spectral cues} and \emph{environmental reverberation}.
Direction-dependent spectral filtering produced by the torso, head, and pinnae introduces characteristic notches and peaks, enabling inference of elevation and disambiguation of positions on the cone of confusion~\cite{gardner19983}.
In addition, the room impulse response (RIR) shaped by early reflections, late reverberation, and surface materials encodes information about source distance, enclosure geometry, and environmental structure~\cite{kuttruff2016room,blauert1997spatial}.
These monaural and environmental cues complement binaural information, jointly supporting robust spatial inference under naturalistic conditions.

Computational models formalize these perceptual cues through structured acoustic representations.
Head-related transfer functions (HRTFs)~\cite{wightman1989headphone} capture listener-specific spectral filtering and are foundational to perceptual research and spatial audio rendering.
In parallel, spherical-harmonics–based ambisonics representations~\cite{zotter2019ambisonics} provide a scene-centric encoding of the sound field around a point in space, with controllable angular resolution.
Together, these representations describe how sound interacts with the body and the environment, enabling principled approaches to spatial reasoning, source localization, and scene reconstruction.
Taken together, the ability of the auditory system to integrate binaural, monaural, and environmental cues demonstrates that sound carries rich implicit structure about the spatial configuration of sources and the surrounding environment.



\begin{figure}[!ht]
    \centering
    \includegraphics[width=\linewidth]{figures/part2/audio_representation.pdf}
    \caption{Three families of acoustic spatial cues: binaural cues (ITD, ILD), monaural spectral cues (direction-dependent filtering by torso, head, and pinnae, summarized by the HRTF), and environmental reverberation (room impulse response).}
    \label{fig:p2_s4_4_audio-representation}
\end{figure}

Recent learning-based approaches aim to capture spatial structure directly from raw audio, going beyond analytic models.
Neural encoders trained on spatially annotated or simulated data learn to infer direction, distance, and environmental geometry from multi-channel waveforms or spectrograms~\cite{chen2022soundspaces,yuksel2025general}.
Learning frameworks further exploit the relation between sound and space: SoundSpaces~\cite{chen2020soundspaces} provides audio-visual simulation for embodied navigation, while VisualEchoes~\cite{gao2020visualechoes} uses image--echo correspondences for self-supervised spatial representation learning.
These learned representations encode not only spatial acoustics but also grounding in surrounding geometry, forming a key component of spatially aware representation.


In real-world environments, spatial audio interacts closely with other sensory modalities.
Audio–visual models align acoustic cues with visual observations to localize sounding objects and understand where events occur~\cite{morgado2018self,afouras2020self,morgado2020learning}.
Embodied agents further leverage auditory signals to perceive layout and navigate toward acoustic targets, as demonstrated in audio–visual navigation~\cite{chen2020soundspaces} and egocentric audio reasoning~\cite{zheng2024bat,biswas2025owl}.
Through such cross-modal grounding, audio provides complementary spatial cues that enhance robustness under visual occlusion and enable reasoning about events beyond the agent's field of view.
This makes audio a powerful complementary modality for spatial representation, especially in scenarios where visual information is partial, noisy, or unavailable.












\section{Domain-Specific Representations}
\label{sec:rep_other}

\FramedInitial{T}{he} representations surveyed so far are organized by modality, each general enough to span many applications. 
A large body of practice, however, develops representations tailored to a specific sensor or downstream task rather than to a modality as such, and these are collected here. 
What unites them is not a shared signal but a shared design logic: each encodes the inductive biases of its data and trades spatial fidelity against semantic compressibility in the way its sensors and target tasks demand. 
This section walks through five such domains: embodied robotics, virtual reality, geospatial and remote sensing, clinical and scientific imaging, and indoor intelligent assistants, examining in each the representations that have become canonical and the constraints that shaped them. 
Beginning with robot perception and planning, systems commonly take as input RGB-D sequences and point clouds (often fused over time into geometric maps that ground observations into a common 3D frame~\cite{newcombe2011kinectfusion, dai2017bundlefusion}), and enable tracking, map building, and manipulation. 
Many such systems also encode visual input into compressed spatial representations, such as keypoints, dense correspondences~\cite{florencemanuelli2018dense}, and point tracks, used to parameterize 6-DoF grasps~\cite{mousavian20196} and object-centric actions.
When available, free space is typically represented with occupancy grids~\cite{Moravec1985SonarMaps, hornung2013octomap} and signed-distance representations~\cite{curless1996volumetric}, and is often additionally augmented with semantic labels such as traversability and cost for planning. 
To further enable action from perception, planners frequently operate on affordance representations~\cite{mo2021where2act} over parts and candidate contact points, encoding feasible interactions such as graspability or predicted contact quality, thereby enabling policies to execute tasks successfully.

Real-time virtual reality systems often represent scenes with triangle meshes and physically based rendering (PBR) materials, where geometry, normals, and texture define appearance, while 6-DoF camera poses keep virtual content registered to the environment. 
To enable flexible interaction, such systems include captured geometry (\eg, spatial meshes, planes, or TSDF and SDF volumes~\cite{curless1996volumetric, newcombe2011kinectfusion}) and depth-based occlusion layers, which are often paired with collision proxies and semantic labels. 
Toward photoreal view synthesis, methods increasingly use neural scene representations (\eg, NeRF-style radiance fields and 3D Gaussian splats~\cite{mildenhall2021nerf, kerbl20233d}), to support novel-view rendering and relighting. 
Additionally, for the purposes of enabling characters in virtual reality systems, such characters are commonly represented with skeletal pose and expression parameters (\eg, SMPL-style body models~\cite{loper2015smpl} or Blendshapes), and their navigation relies on navmeshes combined with acceleration structures (\eg, octrees), to handle moving agents in simulation at high frame rates.

In remote sensing applications, geographic information systems typically represent the world (the Earth) as georeferenced raster fields, for example, optical multispectral and hyperspectral imagery as well as SAR backscatter. 
These fields are aligned to explicit coordinate reference systems and sensor geometry metadata. 
Such systems encode terrain through multiple elevation products (\eg, DEM, DTM, DSM), and compress human infrastructure into vector networks (\eg, roads, power grids, pipelines), often across scales. 
And increasingly, they use LiDAR to not only output raw point clouds, but also to output rasterized height maps that support flexible 3D reasoning. 
These frameworks also commonly operate on multi-temporal stacks and spatiotemporal data with validity masks, which can generate change maps for land-use dynamics and disaster response, as well as feed digital-twin pipelines that generate orthomosaics and photogrammetric point clouds~\cite{schonberger2016structure}, for simulation and virtual modeling.

Clinical and scientific imaging operate on volumetric scans (\eg, CT and MRI), in which 3D voxel grids capture geometry and intensities~\cite{cciccek20163d}, and explicit spatial transforms (\eg, rigid, affine, and nonlinear deformation fields) align subjects and modalities into a shared atlas space. 
In addition, functional measurements (\eg, fMRI BOLD and calcium imaging) are often stored as 4D spatiotemporal tensors, but are frequently compressed into time series and activation maps after parcellation. 
When available, such systems represent structural delineations (\eg, tissue, organ, lesion) as segmentation masks and labeled maps~\cite{ronneberger2015unet}, which are often paired with surface meshes (\eg, cortical surfaces) that enable surface-based topological comparison across subjects. 
For diffusion imaging applications, tensor and orientation fields encode local diffusion directions and magnitudes, while tractography streamlines use these fields to reconstruct fiber pathways and uncover possible structural connectivity. 
These diverse representations across granularities and temporal scales are often mapped onto one another through atlas spaces and parcellations, enabling multimodal fusion and cross-subject analysis.

Intelligent assistants that operate in indoor environments typically ground language in 3D space using multi-scale maps, for example, 3D reconstructions and floorplans anchored by an estimated 6-DoF agent pose~\cite{savva2019habitat, straub2019replica, chang2017matterport3d}, which are often combined with semantic maps that encode navigable routes. 
Context for interactions is captured in structured world models (\eg, scene graphs with typed entities and relations~\cite{armeni20193dsg, wald2020learning, wu2021scenegraphfusion}), which can be augmented with additional semantics like anchors, keypoints, and regions of interest. 
For the purposes of intelligent and persistent assistance, representations that encode some form of memory are necessary, for example, episodic records and long-horizon object trajectories, which are often represented as probabilistic beliefs over object locations in a room. 
In addition, assistant systems can optionally contain geofences and occupancy masks, to help define actionable regions for planning. 

In summary, many domains have developed representations specifically designed for specialized sensors and encoding the inductive biases of the underlying data structure. 
We see that these representations model a tradeoff between (spatial) fidelity and (semantic) compressibility, reflecting how the aforementioned sensors sample space, what priors are built in, and what downstream tasks are to be supported. 



\chapter{Cross-Modal Spatial Representation}
\label{chapter:cross_modality}


% \lettrine[lines=3]{\initfamily\textcolor{darkgreen}{T}}{he} 
\FramedInitial{T}{he}
aim of cross-modal spatial representation is a latent space in which heterogeneous perceptual signals, images, 3D scenes, video, audio, and language are aligned so that the spatial structure each carries is preserved rather than averaged away.
Where the previous chapter asked how much geometry a single modality holds on its own, the concern here is the opposite and harder one: how signals that encode space in incompatible formats (a pixel grid, a world-frame point set, a waveform, a sentence) can be brought into correspondence without the alignment itself distorting what it aligns.

No single recipe achieves this.
Methods differ in what they treat as the unit of alignment, in the objective that binds one modality to another, in how tightly they anchor the space to physical geometry, and in the linguistic interface through which it is queried.
These choices are not independent, and the chapter first makes them explicit as a small set of design axes before turning to the individual cross-modal pairings in which they play out.

\section{Overview}
\label{sec:cross_modal_overview}

\FramedInitial{W}{e} define four axes that structure this cross-modal representation learning problem, which we use to organize the remainder of the chapter. 
The first axis is the \textit{alignment strategy}, which we categorize into contrastive alignment, generative alignment, and programmatic alignment. 
Contrastive alignment learns a shared embedding space by pushing together paired samples and pulling away negative samples. 
Such alignment scales to weakly paired data and enables retrieval and coarse grounding, but often needs additional spatial supervision to understand fine-grained relations grounded to specific visual regions. 
In contrast, generative alignment predicts across modalities (\eg, captioning from vision or language-conditioned decoding into images, 3D, or video), which gives rich supervision and coverage, but also tends to drift from true underlying geometry unless it is further constrained, for example, by multi-view consistency. 
Finally, programmatic alignment aligns visual entities to explicit symbols (\eg, attributes, relation graphs, or concepts in domain-specific languages), which define spatial semantics and naturally interface with program-based, verifiable reasoning, but often requires stronger supervision and trades off expressiveness for structure.

The second axis is \textit{spatial modality}, which ranges from view-aligned 2D to world-aligned 4D. Here, two-dimensional tokens (\eg, pixels, patches) are expressive and broadly applicable, but rely on some form of geometric structure to capture faithful spatial relationships. 
Three-dimensional tokens (\eg, points, meshes, implicit fields, Gaussian primitives) better encode metric scale and rigid-body invariance (\ie SE(3) consistency: rotations and translations of the scene do not change the underlying semantics), so that linguistic queries are grounded in world coordinates rather than in image projections. 
However, they are less ubiquitous and often rely on partial and noisy reconstructions of geometry. 
Extending 3D through time yields 4D traces (\eg, tracks), that bind actions to expressive motion and change, but significantly increase computational complexity and require reliable alignment data, while being sensitive to tracking errors such as identity switches.
These different spatial modalities correspond to and build upon the single-modality representations of Chapter~\ref{chapter:single_modality}.

The third axis is the \textit{language interface}, from free-form text to structured interfaces. 
Free-form text (\eg, natural language) enables open-vocabulary understanding for referring expressions, retrieval, and instruction following, but comes with corresponding ambiguity in terms of frames of reference, scale, and units. 
In contrast, structured interfaces (\eg, programs) reduce such ambiguity by explicitly describing spatial semantics, enabling integration with symbolic reasoning, but are less flexible in capturing the full range of natural language descriptions.

The fourth axis is the \textit{training regime}, which defines both the level of supervision and coverage. Paired supervision, for example, from captioned videos or referring expressions in 3D scenes, gives strong alignment, but is expensive and can be biased toward frequent compositions and co-occurrences. 
Weakly paired internet data, for example, from alt-text or transcripts, enable diversity, but, due to low signal-to-noise ratios, often yield noisier spatial grounding.

Below, we organize the chapter by cross-modal pair, using \emph{spatial modality} (which pair) and \emph{alignment strategy} (which objective) as the primary organisational axes; the \emph{language interface} and \emph{training regime} axes recur as cross-cutting attributes referenced throughout. 
We give an overview of 2D-text alignment (Section~\ref{cross_modality_2d_text}), 3D-text alignment (Section~\ref{cross_modality_3d_text}), audio-text alignment (Section~\ref{cross_modality_audio_text}), and vision-vision alignment (Section~\ref{cross_modality_vision_vision}); each follows the scope and axes defined here.


\section{2D-text}
\label{cross_modality_2d_text}
\FramedInitial{T}{he} most developed cross-modal setting pairs images with language. 
Its defining tension is that a text-aligned 2D representation lives in the image projection, not the world: appearance and layout sit on the pixel grid, yet the relations a caption asserts (left of, above, behind) are only implicit there. 
How faithfully those relations survive alignment turns on two choices: the granularity at which the image is represented, and the objective that binds it to text, and this section is organized around them. 
Models often fall into three broad families of visual representations, across different levels of granularity. First are global embeddings that compress an entire image into a single vector paired with a sentence embedding. 
Here, dual-encoder approaches such as CLIP~\cite{radford2021learning} and ALIGN~\cite{jia2021scaling} show that large-scale contrastive pretraining over noisy web data yields strong zero-shot classification and retrieval performance. 
Second, many works have proposed local tokens tied to image subregions, in which vision-language models, often transformers (\eg, ViLBERT~\cite{lu2019vilbert}, LXMERT~\cite{tan2019lxmert}, UNITER~\cite{chen2020uniter}, OSCAR~\cite{li2020oscar}) input region features extracted by bottom-up detectors (\eg, Faster R-CNN~\cite{ren2015faster}) into cross-modal attention blocks, enabling word-region alignment, and hence fine-grained reasoning over object proposals. 
A third class of methods makes relations explicit; here, models extract intermediate representations that are directly aligned with spatial predicates such as ``left-of'' or ``on.'' NSCL~\cite{mao2019neuro} and LEFT~\cite{hsu2023s} learn hybrid representations that explicitly bind named visual concepts with latent features. 
This representational spectrum, from global to region to relation-based, forms the basis for the objectives below.

Contrastive retrieval serves as a common objective for large-scale pretraining, as it learns a shared metric space, without requiring dense and expensive supervision. 
As representative examples, CLIP~\cite{radford2021learning} and ALIGN~\cite{jia2021scaling} maximize agreement between paired image and text embeddings, enabling robust zero-shot transfer; LiT~\cite{zhai2022lit} shows that contrastively tuning only the text encoder on top of a locked image encoder can further improve such transfer; and CoCa~\cite{yu2022coca} fuses contrastive learning with captioning in a single encoder-decoder architecture to combine discriminative and generative strengths. 
Another line of work adapts contrastive alignment to more fine-grained regions. 
RegionCLIP~\cite{zhong2022regionclip} lifts global contrast to region-text pairs, FILIP~\cite{yao2021filip} proposes late token-token interaction to align words with patches efficiently, X-VLM~\cite{zeng2021multi} learns multi-grained alignment such that entities and attributes can be localized and described consistently across granularities, and ALBEF~\cite{li2021align} conducts image-text contrast with cross-modal fusion and momentum-distilled pseudo-targets to stabilize learning on noisy web data. 

In addition, masked modeling with textual conditioning has been shown to provide a complementary self-supervised signal that encourages cross-modal consistency. LXMERT~\cite{tan2019lxmert} and UNITER~\cite{chen2020uniter} pretrain with masked language modeling, masked region modeling, and image-text matching to build word-region alignment.
VLMo~\cite{bao2022vlmo} generalizes this idea with mixture-of-modality experts, and BEiT-3~\cite{wang2023image} performs masked data modeling over images, text, and image-text pairs in a single multiway transformer.

With more labels, generative objectives can be applied for this representation learning problem, and also complement self-supervised ones. 
Here, methods teach models to produce language conditioned on images, intuitively, to generate tokens that act as spatially grounded explanations. 
As representative examples, SimVLM~\cite{wang2021simvlm} shows that a prefix language-modeling objective trained on weak supervision can improve performance on captioning and visual question answering, BLIP~\cite{li2022blip} adds caption bootstrapping to denoise web supervision, and GIT~\cite{wang2022git} simplifies a scalable image-encoder and text-decoder architecture.

In a more fine-grained manner, grounding objectives can directly optimize spatial correspondence between language and 2D pixels, points, or boxes. Here, text-conditioned detection and phrase grounding unify detection and language. 
MDETR~\cite{kamath2021mdetr} learns a DETR-style~\cite{carion2020end} detector modulated by text, GLIP~\cite{li2022grounded} pretrains on both detection and grounding corpora and self-generates boxes from image-text pairs to scale supervision, and Grounding DINO~\cite{liu2024grounding} augments strong detection backbones with grounded pretraining for open-set detection. 
At the pixel level, DenseCLIP~\cite{rao2022denseclip} adapts image-text matching to pixel-text matching for dense prediction, and GroupViT~\cite{xu2022groupvit} shows that segmentation structure can appear from text supervision alone.

Across these categories, we see that 2D-text representations can emerge from diverse approaches, for example, differing in whether spatial structure is implicit (\eg, in attention over patches) or explicit (\eg, in region tokens), and whether supervision is weakly paired and contrastive, generative, or explicitly grounded. 
Models that learn 2D vision-language representations now increasingly combine all contrastive, captioning, grounding, and masked objectives to solve a wide range of downstream tasks, from (more semantic) referring expression comprehension to (more visually grounded) localization.

\section{3D-text}
\label{cross_modality_3d_text}
\FramedInitial{M}{oving} from the image plane to the scene itself removes the projection gap that limits 2D-text alignment: language is now bound to geometry expressed in world coordinates, so that metric scale and the invariance of meaning under rotation and translation become properties the representation can actually enforce. 
The difficulty shifts accordingly, from recovering relations that pixels hide to obtaining and aligning 3D structure that is often partial, noisy, and expensive to annotate. 
As before, the representation families and the alignment objectives acting on them jointly organize the section. 
Two complementary families of representations have emerged in this domain. First, object-centric and scene-centric encoders learn world-aligned tokens, then align them to language via contrastive or generative objectives. 
ULIP~\cite{xue2023ulip} instantiates this paradigm by aligning point-cloud features to an image-text space, scaling supervision through synthetic triplets.
OpenShape~\cite{liu2023openshape} extends this direction further with curated 3D corpora and hard negative mining to support open-world 3D retrieval, and OpenScene~\cite{peng2023openscene} co-embeds dense features for 3D scenes with text and image pixels in CLIP feature space. 
Additional projection strategies have also been explored, for example, adapting 2D vision-language models to 3D by rendering depth or multi-view images and transferring CLIP's alignment to point clouds; PointCLIP~\cite{zhang2022pointclip} and CLIP2Point~\cite{huang2023clip2point} do so and enable adapters for cross-modal transfer.

A second family binds language to implicit scene fields. Here, LERF~\cite{kerr2023lerf} optimizes a dense 3D language field inside NeRF~\cite{mildenhall2021nerf} by rendering CLIP embeddings across views, enabling open-vocabulary queries that return relevancy volumes. 
CLIP-Fields~\cite{shafiullah2022clip} similarly distills 2D vision-language cues into 3D implicit maps to support semantic search and view localization. 
In addition, replacing NeRF with 3D Gaussian Splatting~\cite{kerbl20233d} yields efficient language fields; for example, LangSplat~\cite{qin2024langsplat} combines CLIP features with Gaussians for fast, precise open-vocabulary querying. 
These methods build on advances in explicit 3D primitives and real-time rendering to maintain high-fidelity representations while training efficiently.

In this 3D-text domain, objectives also follow the representation choices. Contrastive alignment remains common, for example, ULIP trains with InfoNCE~\cite{oord2018representation} over point clouds-images-text, while OpenShape scales both data and hard-negative mining. 
Projection-based approaches transfer contrastive objectives to depth and multi-view renders to align point clouds with text labels. 
Grounded objectives enforce explicit 3D correspondence between language and geometry. 
3D referring expression datasets and models (\eg, ScanRefer~\cite{chen2020scanrefer}, ReferIt3D~\cite{achlioptas2020referit3d}, 3D-VisTA~\cite{zhu20233d}, BUTD-DETR~\cite{jain2022bottom}, NS3D~\cite{hsu2023ns3d}, InstanceRefer~\cite{yuan2021instancerefer}) localize targets from language descriptions and learn to understand complex relations dependent on viewpoints.

Here, geometry-aware representations are critical for aligning language with the physical world. 
Ideal 3D-text representations should respect SE(3) structure (\eg, rotations and translations should not spuriously change semantics), yet cross-modal alignment should still encode view-dependent language (\eg, ``the chair to the left of the table''). 
To this end, priors from equivariant architectures, such as the SE(3)-Transformer~\cite{fuchs2020se} and Vector Neurons~\cite{deng2021vector}, can be integrated into 3D-text models. 
Implicit fields enable different options; while NeRF-based language fields are not strictly equivariant, multi-view consistency and cross-view regularization can improve their spatial coherence. 
Alternatively, Gaussian splatting uses a canonical world frame and enables efficient supervision during language-field optimization. 
Overall, these representation choices and objectives give rise to a wide range of approaches for 3D-text alignment, from tri-modal contrastive pretraining to implicit or splat-based language fields.

\section{Audio-text}
\label{cross_modality_audio_text}
\FramedInitial{S}{ound} aligns to language along a different primary dimension than vision does. What an acoustic signal encodes first is time, the sequence of events as they unfold, and only secondarily, and coarsely, does it carry spatial attributes such as direction of arrival or distance. 
Audio-text representation learning therefore aims for embeddings that are open-vocabulary and temporally grounded, with spatial awareness added where the labels permit. 
The section again moves from the representation granularity to the objectives that align it with text. 
One such direction is to embed audio-text representations at a coarse level, for example, with global encoders mapping an entire audio clip and a full sentence into a joint space, enabling zero-shot retrieval. 
Here, CLAP~\cite{elizalde2023clap} and LAION-CLAP~\cite{wu2023large} scale contrastive language-audio pretraining with curated datasets, and show strong transfer across retrieval benchmarks. 
In addition, AudioCLIP~\cite{guzhov2022audioclip} extends CLIP to a tri-modal image-text-audio setting, to leverage visual-text semantics, and Wav2CLIP~\cite{wu2022wav2clip} distills CLIP's alignment into an audio encoder when audio-text pairs are scarce.

Beyond global representations, a line of work that aims to extract representations at a more fine-grained scale learns temporally local structure, for example, by operating on spectrogram patch tokens with transformer architectures. 
The Audio Spectrogram Transformer (AST)~\cite{gong2021ast} and PaSST~\cite{koutini2021efficient} show that attention over time-frequency patches captures long-range temporal dependencies, and HTS-AT~\cite{chen2022hts} further returns frame-level maps that are useful for event localization. 
These tokenizations make it straightforward to bind subclip semantics to text (\eg, ``a door slams, then footsteps echo''). 
Self-supervised masked objectives, \eg, SSAST~\cite{gong2022ssast}, MAE-AST~\cite{baade2022mae}, and Audio-MAE~\cite{huang2022masked}, additionally pretrain such backbones by reconstructing masked patches, while BEATs~\cite{chen2022beats} adds an iterative tokenizer-encoder scheme that encourages high-level semantic abstraction.

Across advances in representations and objectives, an emergent and increasingly common paradigm is to first pretrain an audio backbone with masked modeling, then align embeddings to text via contrastive losses, and finally, optionally add a captioning head to increase supervision. 
In tasks where spatial semantics is important, incorporating spatial labels, such as direction-of-arrival, helps refine the audio-text representation space. 
By leveraging these approaches, models can learn to align predicates such as ``left-of'' or ``far'' with meaningful acoustic features. 
These methods enable models to learn embeddings that are not only open-vocabulary and temporally grounded, but also (coarsely) spatially aware.

\section{Vision-vision}
\label{cross_modality_vision_vision}
\FramedInitial{A}{} common approach to improving vision-language alignment is to first align vision with vision, intuitively, enforcing geometric consistency such that visual features are anchored to a persistent world model before combining with language.
Concretely, the goal of vision-vision representation learning is to fuse visual inputs (\eg, multiview RGB) into a set of world-aligned scene tokens, such that they are invariant to viewpoint changes and can be later queried by language. 
Here, we discuss three broad families of approaches: lifting 2D features into 3D scene tokens, projecting these 3D tokens back into 2D for supervision, and using audiovisual correspondence to add constraints that further anchor the world model.

Lifting 2D to 3D follows two complementary directions. The most classical and canonical approach first estimates camera poses and geometry with structure-from-motion (SfM) and simultaneous localization and mapping (SLAM), then adds semantics to the generated map. 
Here, commonly used frameworks include COLMAP for SfM~\cite{schonberger2016structure}, KinectFusion and ElasticFusion for dense RGB-D mapping~\cite{newcombe2011kinectfusion,  whelan2015elasticfusion}, and ORB-SLAM2 for SLAM~\cite{mur2017orb}. 
Learning-based reconstructions build upon this direction by back-projecting per-view features into a 3D volume and predicting TSDF or occupancy directly from posed images, for example, as in NeuralRecon~\cite{sun2021neuralrecon}. 
Though components may differ, all of these approaches output view-agnostic scene tokens that can later be guided by linguistic supervision and queries.

Projecting 3D back to 2D provides self-supervision. Differentiable rendering in neural fields anchors world-aligned tokens to multiview observations, and enforces cross-view consistency~\cite{mildenhall2021nerf}. 
Distilling 2D foundation features into 3D fields adds semantic knowledge without manual, expensive 3D labels. Here, works typically transfer 2D features into a co-registered 3D feature volume, which can be rendered to any view~\cite{shen2023distilled, ye2023featurenerf}. 
In addition, distillation is not limited to 2D language targets alone; for example, ImageBind~\cite{girdhar2023imagebind} proposes a multisensor feature target (image-text-audio-depth-thermal-IMU), which enables transfer of cross-sensor semantics even when there are few explicit labels. 
Although these distillations are vision-vision processes, they are explicitly designed for language alignment later. The key benefit is that now vision-language alignment is conducted on a geometry-consistent representation, rather than directly on view-specific features.

Audio-visual alignment can further constrain this vision-vision world model; self-supervised objectives have been designed to enforce consistency between audio and visual motion, which allows models to localize sound sources and disambiguate co-occurring events~\cite{arandjelovic2017look, owens2018audio, zhao2018sound}. 
Such cues provide directional or timing semantics that help anchor linguistic queries to specific world-aligned entities, especially in cases with multiple candidate entities. 
Across this variety of approaches, vision-vision representation learning methods all aim to learn a persistent world model that is geometrically and temporally consistent before adding any linguistic supervision, therefore making downstream cross-modal alignment more data-efficient and spatially faithful to the physical world.


